\documentclass[11pt,letterpaper]{article}
\usepackage[top=0.7in,bottom=0.7in,left=0.75in,right=0.75in]{geometry}
\usepackage[T1]{fontenc}
\usepackage[utf8]{inputenc}
\usepackage{lmodern}
\usepackage{microtype}
\usepackage{enumitem}
\usepackage[hidelinks]{hyperref}
\usepackage{titlesec}
\setlength{\parindent}{0pt}
\setlength{\parskip}{0pt}
\pagestyle{empty}
\titleformat{\section}{\normalfont\bfseries}{}{0em}{\uppercase}[\vspace{-4pt}\hrule\vspace{4pt}]
\titlespacing{\section}{0pt}{8pt}{5pt}
\setlist[itemize]{leftmargin=1.2em,label=--,itemsep=0pt,topsep=2pt,parsep=0pt}
\begin{document}
\begin{center}
{\fontsize{16}{19}\selectfont\textbf{\MakeUppercase{Diego Castillo Pineda}}}\\[4pt]
{\small La Florida, Santiago, Chile \quad|\quad +56 9 9416 2547 \quad|\quad \href{mailto:diego.castillop11@gmail.com}{diego.castillop11@gmail.com}\\
\href{https://https://www.linkedin.com/in/diego-castillo-pineda/}{https://www.linkedin.com/in/diego-castillo-pineda/} \quad|\quad \href{https://https://github.com/diegocastillop11-cloud}{https://github.com/diegocastillop11-cloud}}
\end{center}

\section{Resumen Profesional}
Analista de Datos especializado en SQL Server, Python y automatización de procesos para Supply Chain, con trayectoria en entornos de alta concurrencia y manejo de grandes volúmenes transaccionales. Experiencia comprobada en diseño de stored procedures que procesaron decenas de miles de transacciones por hora, optimización de consultas con SQL Server Management Studio, y desarrollo de reportería avanzada con Power BI y Excel. Capacidad para integrar datos entre sistemas ERP (Oracle, SAP) y plataformas cloud (Microsoft Azure), aplicando Business Analytics y fundamentos de Data Science para respaldar decisiones operativas en logística y cadena de suministro.

\section{Experiencia Profesional}

\textbf{Punto Ticket} \hfill Santiago, Chile\\
\textit{Analista de Base de Datos Senior} \hfill Nov. 2019 – Feb. 2026
\begin{itemize}
  \item Diseñé y optimicé stored procedures en SQL Server para eventos masivos, procesando decenas de miles de transacciones por hora y manteniendo tiempos de respuesta bajo 2 segundos durante picos de carga.
  \item Automaticé procesos masivos mediante scripts T-SQL, reduciendo errores manuales en un 85\% y mejorando la eficiencia operativa en operaciones recurrentes de alta concurrencia.
  \item Desarrollé consultas dinámicas con PIVOT y reportería avanzada en Power BI, reemplazando procesos manuales y entregando insights de negocio en tiempo real para la toma de decisiones estratégicas.
  \item Integré datos entre sistemas heterogéneos utilizando HTML, JSON y XML, asegurando la disponibilidad y consistencia de la información en entornos productivos críticos.
  \item Participé en migración de bases de datos hacia Microsoft Azure, colaborando con equipos técnicos para garantizar la continuidad operativa y la adopción de arquitecturas cloud.
  \item Administré bases de datos SQL Server en entornos de alta disponibilidad, implementando mejores prácticas de seguridad, respaldo y monitoreo continuo durante períodos críticos de operación.
\end{itemize}
\vspace{2pt}
\textbf{Imperial S.A.} \hfill Santiago, Chile\\
\textit{Analista Funcional} \hfill Jun. 2017 – Nov. 2019
\begin{itemize}
  \item Realicé análisis funcional y soporte a distintas áreas de negocio utilizando SQL Server Management Studio, Oracle y Visual Basic 6, resolviendo incidencias que impactaban directamente la operación productiva.
  \item Levanté requerimientos junto a usuarios clave, parametrizando y configurando sistemas ERP según especificaciones funcionales, lo que mejoró la trazabilidad de procesos logísticos internos.
  \item Elaboré documentación técnica detallada (instructivos, flujos de proceso, manuales de usuario) y capacité a usuarios finales, facilitando la adopción de nuevas funcionalidades y reduciendo tiempos de respuesta en soporte.
  \item Colaboré con equipos técnicos y funcionales en reuniones de seguimiento, asegurando la correcta implementación de cambios y la alineación entre TI y áreas de negocio en proyectos de mejora continua.
\end{itemize}
\vspace{2pt}
\textbf{OB GROUP PARK} \hfill Santiago, Chile\\
\textit{Analista de Sistemas y TI} \hfill Mayo 2014 – Abril 2017
\begin{itemize}
  \item Lideré proyectos de mejora en procesos contables y de gestión retail, coordinando la integración de nuevas tecnologías y asegurando la compatibilidad entre plataformas ERP y sistemas legacy.
  \item Colaboré con áreas de finanzas y operaciones para garantizar la continuidad operativa durante implementaciones tecnológicas, minimizando interrupciones y facilitando la adopción de cambios.
  \item Desarrollé reportes y consultas SQL personalizadas para áreas de negocio, mejorando la visibilidad sobre inventarios y movimientos logísticos en puntos de venta.
\end{itemize}
\vspace{2pt}
\textbf{Hewlett-Packard Company} \hfill Santiago, Chile\\
\textit{Agente de Mesa de Ayuda Nivel 1} \hfill Ago. 2013 – Mayo 2014
\begin{itemize}
  \item Brindé soporte técnico a usuarios finales, desarrollando habilidades de diagnóstico y resolución de problemas que fueron la base de mi trayectoria profesional en TI.
\end{itemize}

\section{Proyectos Destacados}

\textbf{Sistema de Gestión Inteligente de Inventarios con IA} \hfill 2024
\begin{itemize}
  \item Desarrollé una aplicación web con React y Node.js que integra Claude API para generar recomendaciones automáticas de reabastecimiento en cadenas de suministro, utilizando PostgreSQL como base de datos segura.
  \item Implementé prompt engineering avanzado con Claude y Gemini para analizar patrones de demanda histórica y predecir necesidades de inventario, mejorando la precisión de pronósticos en un 30\% frente a métodos tradicionales.
  \item Diseñé una interfaz UI/UX intuitiva con TypeScript y visualización de datos en tiempo real, permitiendo a usuarios no técnicos consultar y validar recomendaciones generadas por IA.
  \item Automaticé el mantenimiento de datos mediante scripts Python que limpian, validan y actualizan registros en PostgreSQL, asegurando la integridad y calidad de la información para análisis predictivos.
  \item Desplegué la solución en Docker y Microsoft Azure, garantizando escalabilidad y disponibilidad continua para pruebas de concepto en entornos de logística y Supply Chain.
\end{itemize}

\section{Habilidades T\'{e}cnicas}
\textbf{Datos \& Supply Chain:} SQL Server (Experto), Azure SQL, Oracle, PostgreSQL, MongoDB, Power BI, Excel Avanzado, Business Analytics, ETL, Data Warehousing \\
\textbf{Análisis \& IA:} Python (Pandas, NumPy), Machine Learning, Prompt Engineering (Claude, Gemini), Data Science, Optimización de Procesos, Forecasting \\
\textbf{Desarrollo \& Cloud:} C\#, Node.js, Go, React, TypeScript, JavaScript (ES6+), Microsoft Azure, Docker, Git (GitHub/GitLab) \\
\textbf{Frontend \& UX:} HTML5/CSS3, UI/UX Design, Visualización de Datos \\
\textbf{Herramientas \& Metodologías:} SQL Server Management Studio, Stored Procedures, T-SQL, JSON/XML, Scrum/Kanban, Documentación Técnica

\section{Formaci\'{o}n Acad\'{e}mica}
\textbf{Business Analytics Essentials, Machine Learning con Python, Data Science y Power BI, Optimización de Procesos y Administración de Bases de Datos SQL Server}, Escuela de Postgrado, Universidad de Chile \hfill 2023 \\
\textbf{Desarrollo de Apps Móviles}, Universidad Complutense de Madrid \hfill 2017 \\
\textbf{Analista Programador}, Universidad Tecnológica de Chile Inacap \hfill 2013

\end{document}